% GENERATED FILE. Edit canonical lessons, then run npm run generate:books.
% canonical-source-hash: fnv1a64:17a00dc2c4c943cf
% canonical-lessons: KA-C534-raga, KA-C534-tala, KA-C534-svara, KA-C534-gunugu, KA-C534-nudisu

\chapter{Music}
\label{ch:ka-music}
\hlchaptermodality{534}

\begin{chapteropening}
\textbf{By the end of this chapter:} \emph{I can say a melody, a rhythm, a musical note, to hum a tune and to play an instrument.}

Complete the last lesson of chapter 534: I can say a melody, a rhythm, a musical note, to hum a tune and to play an instrument.
\end{chapteropening}

\section[\texorpdfstring{\kn{ರಾಗ} (rāga)}{ರಾಗ (rāga)}]{\kn{ರಾಗ} (rāga) — a melody; a raga}

\label{lesson:KA-C534-raga}



\begin{quote}
Before the new one: say the Kannada for to string beads, then the Kannada for to adorn.
\end{quote}

\subsection*{You'll want to know: \kn{ರಾಗ}}
\textbf{\kn{ರಾಗ}} — \emph{rāga} — \textquotedblleft{}a melody; a raga\textquotedblright{}.

\begin{quote}
\kn{ಆ} \kn{ಗಾಯಕ} \kn{ಸಂಜೆಯ} \kn{ರಾಗವನ್ನು} \kn{ಹಾಡಿದರು}.
\end{quote}

\begin{quote}
\emph{ā gāyaka sañjeya rāgavannu hāḍidaru.}
\end{quote}

\begin{quote}
That singer sang an evening raga.
\end{quote}

Each \kn{ರಾಗ} has its own set of notes and mood.

\begin{cousinweb}
\textbf{\kn{ರಾಗ}} (\emph{rāga}) is borrowed from Sanskrit \emph{rāga}, \textquotedblleft{}colour\textquotedblright{}, and \textquotedblleft{}passion, love\textquotedblright{} (CDIAL 10674): a melody that colours the mind.
\end{cousinweb}

\begin{grammarlens}[title={What you've built}]
One more word for describing what you lived and felt.
\end{grammarlens}

\subsection*{Your turn}
\begin{itemize}
\raggedright
  \item \emph{Say it:} \emph{rāga}
  \item \emph{Say it:} \emph{rāga}, once more
  \item \emph{Say it:} \emph{ā gāyaka sañjeya rāgavannu hāḍidaru.} — That singer sang an evening raga.
  \item \emph{Recall:} say the Kannada for to string beads, then the Kannada for to adorn, then say \emph{rāga} again
\end{itemize}

\begin{tcolorbox}[breakable,colback=teal!4,colframe=teal!35!black,title={Before you move on}]
What does \kn{ರಾಗ} mean? (\textquotedblleft{}A melody; a raga\textquotedblright{}.) Say it once more.
\end{tcolorbox}
\section[\texorpdfstring{\kn{ತಾಳ} (tāḷa)}{ತಾಳ (tāḷa)}]{\kn{ತಾಳ} (tāḷa) — a rhythm; a beat}

\label{lesson:KA-C534-tala}



\begin{quote}
Before the new one: say the Kannada for to adorn, then the Kannada for a melody.
\end{quote}

\subsection*{You'll want to know: \kn{ತಾಳ}}
\textbf{\kn{ತಾಳ}} — \emph{tāḷa} — \textquotedblleft{}a rhythm; a beat\textquotedblright{}.

\begin{quote}
\kn{ಮಕ್ಕಳು} \kn{ಹಾಡಿಗೆ} \kn{ತಾಳ} \kn{ಹಾಕಿದರು}.
\end{quote}

\begin{quote}
\emph{makkaḷu hāḍige tāḷa hākidaru.}
\end{quote}

\begin{quote}
The children kept time to the song.
\end{quote}

\kn{ತಾಳ} \kn{ಹಾಕು}, to keep time.

\begin{cousinweb}
\textbf{\kn{ತಾಳ}} (\emph{tāḷa}) is borrowed from Sanskrit \emph{tāla}, \textquotedblleft{}a cymbal\textquotedblright{} (CDIAL 5801): the beat the cymbals keep.
\end{cousinweb}

\begin{grammarlens}[title={What you've built}]
One more word for describing what you lived and felt.
\end{grammarlens}

\subsection*{Your turn}
\begin{itemize}
\raggedright
  \item \emph{Say it:} \emph{tāḷa}
  \item \emph{Say it:} \emph{tāḷa}, once more
  \item \emph{Say it:} \emph{makkaḷu hāḍige tāḷa hākidaru.} — The children kept time to the song.
  \item \emph{Recall:} say the Kannada for to adorn, then the Kannada for a melody, then say \emph{tāḷa} again
\end{itemize}

\begin{tcolorbox}[breakable,colback=teal!4,colframe=teal!35!black,title={Before you move on}]
What does \kn{ತಾಳ} mean? (\textquotedblleft{}A rhythm; a beat\textquotedblright{}.) Say it once more.
\end{tcolorbox}
\section[\texorpdfstring{\kn{ಸ್ವರ} (svara)}{ಸ್ವರ (svara)}]{\kn{ಸ್ವರ} (svara) — a musical note; a voice}

\label{lesson:KA-C534-svara}



\begin{quote}
Before the new one: say the Kannada for a melody, then the Kannada for a rhythm.
\end{quote}

\subsection*{You'll want to know: \kn{ಸ್ವರ}}
\textbf{\kn{ಸ್ವರ}} — \emph{svara} — \textquotedblleft{}a musical note; a voice\textquotedblright{}.

\begin{quote}
\kn{ಸಂಗೀತದಲ್ಲಿ} \kn{ಏಳು} \kn{ಸ್ವರಗಳು} \kn{ಇವೆ}.
\end{quote}

\begin{quote}
\emph{saṅgītadalli ēḷu svaragaḷu ive.}
\end{quote}

\begin{quote}
There are seven notes in music.
\end{quote}

Plural: \kn{ಸ್ವರಗಳು}, \emph{svaragaḷu}.

Also a vowel, and the tone of a voice.

\begin{cousinweb}
\textbf{\kn{ಸ್ವರ}} (\emph{svara}) is borrowed from Sanskrit \emph{svara}, \textquotedblleft{}sound; a musical note\textquotedblright{} (CDIAL 13908).
\end{cousinweb}

\begin{grammarlens}[title={What you've built}]
One more word for describing what you lived and felt.
\end{grammarlens}

\subsection*{Your turn}
\begin{itemize}
\raggedright
  \item \emph{Say it:} \emph{svara}
  \item \emph{Say it:} \emph{svara}, once more
  \item \emph{Say it:} \emph{saṅgītadalli ēḷu svaragaḷu ive.} — There are seven notes in music.
  \item \emph{Recall:} say the Kannada for a melody, then the Kannada for a rhythm, then say \emph{svara} again
\end{itemize}

\begin{tcolorbox}[breakable,colback=teal!4,colframe=teal!35!black,title={Before you move on}]
What does \kn{ಸ್ವರ} mean? (\textquotedblleft{}A musical note; a voice\textquotedblright{}.) Say it once more.
\end{tcolorbox}
\section[\texorpdfstring{\kn{ಗುನುಗು} (gunugu)}{ಗುನುಗು (gunugu)}]{\kn{ಗುನುಗು} (gunugu) — to hum a tune}

\label{lesson:KA-C534-gunugu}



\begin{quote}
Before the new one: say the Kannada for a rhythm, then the Kannada for a musical note.
\end{quote}

\subsection*{You'll want to know: \kn{ಗುನುಗು}}
\textbf{\kn{ಗುನುಗು}} — \emph{gunugu} — \textquotedblleft{}to hum a tune\textquotedblright{}.

\begin{quote}
\kn{ಅಡುಗೆ} \kn{ಮಾಡುವಾಗ} \kn{ಅಮ್ಮ} \kn{ಒಂದು} \kn{ಹಾಡು} \kn{ಗುನುಗುತ್ತಿದ್ದರು}.
\end{quote}

\begin{quote}
\emph{aḍuge māḍuvāga amma ondu hāḍu gunuguttiddaru.}
\end{quote}

\begin{quote}
While cooking, mother was humming a song.
\end{quote}

Past (I …): \kn{ಗುನುಗಿದೆ}, \emph{gunugide}.

Softly, with the lips nearly closed.

\begin{cousinweb}
\textbf{\kn{ಗುನುಗು}} (\emph{gunugu}) imitates the sound itself; DEDR 1685 lists the related Kannada \emph{guṇuguṇu}, \textquotedblleft{}murmuring, buzzing\textquotedblright{}. The exact origin is not settled.
\end{cousinweb}

\begin{grammarlens}[title={What you've built}]
One more word for describing what you lived and felt.
\end{grammarlens}

\subsection*{Your turn}
\begin{itemize}
\raggedright
  \item \emph{Say it:} \emph{gunugu}
  \item \emph{Say it:} \emph{gunugu}, once more
  \item \emph{Say it:} \emph{aḍuge māḍuvāga amma ondu hāḍu gunuguttiddaru.} — While cooking, mother was humming a song.
  \item \emph{Recall:} say the Kannada for a rhythm, then the Kannada for a musical note, then say \emph{gunugu} again
\end{itemize}

\begin{tcolorbox}[breakable,colback=teal!4,colframe=teal!35!black,title={Before you move on}]
What does \kn{ಗುನುಗು} mean? (\textquotedblleft{}To hum a tune\textquotedblright{}.) Say it once more.
\end{tcolorbox}
\section[\texorpdfstring{\kn{ನುಡಿಸು} (nuḍisu)}{ನುಡಿಸು (nuḍisu)}]{\kn{ನುಡಿಸು} (nuḍisu) — to play an instrument}

\label{lesson:KA-C534-nudisu}



\begin{quote}
Before the new one: say the Kannada for a musical note, then the Kannada for to hum a tune.
\end{quote}

\subsection*{You'll want to know: \kn{ನುಡಿಸು}}
\textbf{\kn{ನುಡಿಸು}} — \emph{nuḍisu} — \textquotedblleft{}to play an instrument\textquotedblright{}.

\begin{quote}
\kn{ನನ್ನ} \kn{ಅಣ್ಣ} \kn{ಗಿಟಾರ್} \kn{ನುಡಿಸುತ್ತಾನೆ}.
\end{quote}

\begin{quote}
\emph{nanna aṇṇa giṭār nuḍisuttāne.}
\end{quote}

\begin{quote}
My elder brother plays the guitar.
\end{quote}

Past (I …): \kn{ನುಡಿಸಿದೆ}, \emph{nuḍiside}.

Literally to make it speak: \kn{ನುಡಿ}, to speak.

\begin{cousinweb}
\textbf{\kn{ನುಡಿಸು}} (\emph{nuḍisu}) is native Dravidian: DEDR 3784 lists Kannada \emph{nuḍi}, \textquotedblleft{}to sound, utter, speak\textquotedblright{}, and \emph{nuḍisu}, \textquotedblleft{}to cause to utter\textquotedblright{}: you make the instrument speak.
\end{cousinweb}

\begin{grammarlens}[title={What you've built}]
One more word for describing what you lived and felt.
\end{grammarlens}

\subsection*{Your turn}
\begin{itemize}
\raggedright
  \item \emph{Say it:} \emph{nuḍisu}
  \item \emph{Say it:} \emph{nuḍisu}, once more
  \item \emph{Say it:} \emph{nanna aṇṇa giṭār nuḍisuttāne.} — My elder brother plays the guitar.
  \item \emph{Recall:} say the Kannada for a musical note, then the Kannada for to hum a tune, then say \emph{nuḍisu} again
\end{itemize}

\begin{tcolorbox}[breakable,colback=teal!4,colframe=teal!35!black,title={Before you move on}]
What does \kn{ನುಡಿಸು} mean? (\textquotedblleft{}To play an instrument\textquotedblright{}.) Say it once more.
\end{tcolorbox}
